\chapter{Research Method}
\label{chap:research_method}

The purpose of this chapter is to describe the research design and convince the reader that it is
appropriate to answer the research questions and validate the thesis statement. The chapter follows
the structure of the Research Onion, presenting only the choices made for this study.

\begin{table}[htbp]
  \centering
  \caption{Research design choices}
  \label{tab:research_choices}
  \begin{tabular}{@{}ll@{}}
    \toprule
    \textbf{Layer}                & \textbf{Choice} \\
    \midrule
    Philosophy (only 1)           & Pragmatism \\
    Approach (only 1)             & Abductive \\
    Methodological Choice (only 1)& Mono — Quantitative (artifact + evaluation) \\
    Strategy (only 1)             & Design Science Research (DSR) \\
    Time Horizon (only 1)         & Cross-sectional \\
    Data Collection (1 or more)   & Historical staffing data; expert interviews \\
    Data Analysis (1 or more)     & ML model training \& evaluation; thematic coding \\
    \bottomrule
  \end{tabular}
\end{table}

% ─────────────────────────────────────────────────────────────────────────────
\section{Philosophy}
\label{sec:philosophy}

The philosophical stance for this research is \textbf{pragmatism}. Pragmatism is appropriate
because the primary aim is to develop a working artifact (the hybrid AI recommendation engine)
that solves a concrete organizational problem. Truth is assessed by practical outcomes — whether
the system demonstrably improves bottleneck anticipation — rather than by adherence to a single
ontological or epistemological paradigm. This stance allows combining qualitative domain
knowledge elicitation (symbolic AI design) with quantitative model evaluation (subsymbolic AI
training and testing).

% ─────────────────────────────────────────────────────────────────────────────
\section{Approach}
\label{sec:approach}

An \textbf{abductive} approach is followed: starting from observed anomalies in CAPEX resource
planning practice, the research moves iteratively between theory (ontologies, ML forecasting
literature) and empirical data (project staffing records) to build and refine the best
explanatory model. This differs from pure induction (no prior theory) and pure deduction
(testing a fixed hypothesis without iteration), matching the iterative nature of Design Science
Research cycles.

% ─────────────────────────────────────────────────────────────────────────────
\section{Methodological Choice}
\label{sec:methodological_choice}

The methodological choice is \textbf{mono-method with a quantitative emphasis}: the core
research activity is the design, implementation, and quantitative evaluation of an IT artifact.
Qualitative elements (expert interviews for domain knowledge acquisition and requirements
validation) are used instrumentally within the DSR process and do not constitute a separate
research strand.

% ─────────────────────────────────────────────────────────────────────────────
\section{Strategy}
\label{sec:strategy}

The strategy is \textbf{Design Science Research (DSR)} \citep{hevner2004design}, which is
appropriate for the development and evaluation of IT artifacts in applied organizational
contexts. The primary artifact is the hybrid AI recommendation engine; the research contribution
lies in both the artifact itself and the generalizable design knowledge extracted from its
development and evaluation. DSR is structured around three cycles: the \textit{Relevance Cycle}
(problem and requirements from practice), the \textit{Design Cycle} (build and evaluate),
and the \textit{Rigor Cycle} (grounding in existing knowledge base). This maps directly onto
the AI-4-SME Design–Build–Run phases \citep{peter2025ai4sme}.

% ─────────────────────────────────────────────────────────────────────────────
\section{Time Horizon}
\label{sec:time_horizon}

The time horizon is \textbf{cross-sectional}: data collection, model training, and evaluation
are performed within a defined period (approx.\ January 2026 to September 2026). This restricts
the ability to capture long-term adoption trends and reduces generalizability of performance
results; this is explicitly acknowledged as a limitation (see Section~\ref{sec:delineations}).

% ─────────────────────────────────────────────────────────────────────────────
\section{Data Collection}
\label{sec:data_collection}

Two data collection methods are used:

\begin{enumerate}
  \item \textbf{Historical project staffing data:} Anonymized records from a partner
        organization (or synthetic benchmarks generated from industry data) covering project
        metadata, role assignments, skill tags, and actuals vs.\ planned allocations. This
        data underpins the quantitative demand forecasting models.

  \item \textbf{Expert interviews / workshops:} Semi-structured interviews with portfolio
        managers and resource planners to (a)~elicit domain knowledge for ontology design,
        (b)~validate the recommendation engine output against expert judgment, and (c)~assess
        perceived decision support quality. The following guiding questions structure the
        interviews:
\end{enumerate}

\begin{table}[htbp]
  \centering
  \caption{Interview guiding questions}
  \label{tab:interview_questions}
  \begin{tabular}{@{}p{9cm}p{5.5cm}@{}}
    \toprule
    \textbf{Question} & \textbf{Purpose} \\
    \midrule
    How do you currently identify upcoming resource bottlenecks in the portfolio? &
    Understand baseline process and pain points \\[0.4em]
    What data sources do you rely on for staffing decisions, and how reliable are they? &
    Identify available data assets and quality issues \\[0.4em]
    Which skill categories are most difficult to plan for and why? &
    Prioritize ESCO coverage requirements \\[0.4em]
    What would an ideal recommendation look like, and what would make you trust it? &
    Define output requirements and explainability needs \\[0.4em]
    What constraints (budget, regulation, HR policy) must any recommendation respect? &
    Capture constraint requirements for allocation engine \\
    \bottomrule
  \end{tabular}
\end{table}

% ─────────────────────────────────────────────────────────────────────────────
\section{Data Analysis}
\label{sec:data_analysis}

\begin{itemize}
  \item \textbf{Quantitative:} ML model training and evaluation using standard metrics
        (MAE, RMSE for forecasting; precision/recall for bottleneck classification).
        Comparison against a baseline (moving average / current ERP planning output).

  \item \textbf{Qualitative:} Thematic coding of interview transcripts to extract design
        requirements and validation feedback. Codes are organized around the six W-questions
        of the AI-4-SME task design method.
\end{itemize}

% ─────────────────────────────────────────────────────────────────────────────
\section{Methodological Framework: AI-4-SME Design--Build--Run}
\label{sec:ai4sme_phases}

The AI-4-SME Framework \citep{peter2025ai4sme} structures the research across three phases:

\subsection*{Phase 1 — Design}

\begin{itemize}
  \item \textbf{Company level:} Identify CAPEX resource management as the strategic AI
        application domain; assess tasks for AI suitability using the feasibility-vs-impact
        matrix.
  \item \textbf{Process level:} Map the CAPEX project lifecycle; identify KIAs
        (skill-demand forecasting, bottleneck detection) and DIAs (allocation optimization,
        capacity computation).
  \item \textbf{Task level:} Apply the six W-questions to define the point-of-view for each
        AI task; generate solution ideas and select the hybrid path based on effort-vs-benefit.
\end{itemize}

\subsection*{Phase 2 — Build (Hybrid AI)}

\begin{enumerate}
  \item \textbf{Knowledge-based stream:} Acquire and represent domain knowledge using ESCO
        as the skill ontology backbone; implement a knowledge graph (Neo4j or GraphDB).
  \item \textbf{Data-based stream:} Prepare historical staffing data; train and evaluate ML
        models (gradient boosting, LSTM, Temporal Fusion Transformer) for temporal
        skill-demand forecasting using knowledge graph embeddings as features.
  \item \textbf{Hybrid integration:} Combine both streams into a unified recommendation
        engine delivering ontology-grounded, data-driven allocation recommendations.
\end{enumerate}

\subsection*{Phase 3 — Run}

\begin{itemize}
  \item \textbf{Governance:} Evaluate GDPR and EU AI Act compliance (transparency,
        explainability, risk classification).
  \item \textbf{IT integration:} Define API interface for ERP and PPM system integration.
  \item \textbf{Human-AI collaboration:} Design the recommendation interface to support —
        not replace — portfolio manager decision-making.
\end{itemize}

% ─────────────────────────────────────────────────────────────────────────────
\section{Work Plan, Milestones, and Risks}
\label{sec:workplan}

\begin{table}[htbp]
  \centering
  \caption{Work packages and milestones}
  \label{tab:workplan}
  \begin{tabular}{@{}p{0.5cm}p{5.5cm}p{2.5cm}p{2cm}p{3cm}@{}}
    \toprule
    \textbf{\#} & \textbf{Task / Work Package} & \textbf{Output} &
    \textbf{Duration} & \textbf{Milestone} \\
    \midrule
    1 & Literature review \& gap analysis           & Screened corpus, gap statement & 4 weeks  & M1: Literature finalized \\
    2 & Domain knowledge elicitation (interviews)   & Interview data, ontology requirements & 3 weeks & M2: Requirements validated \\
    3 & ESCO knowledge graph design \& build        & Knowledge graph prototype & 4 weeks  & M3: KG operational \\
    4 & Data preparation \& feature engineering     & ML-ready dataset & 3 weeks  & \\
    5 & ML model development \& training            & Trained forecasting models & 5 weeks  & M4: Models benchmarked \\
    6 & Hybrid integration \& recommendation engine & Integrated prototype & 4 weeks  & M5: Prototype complete \\
    7 & Evaluation \& expert validation             & Evaluation results & 3 weeks  & M6: Evaluation done \\
    8 & Thesis writing \& revision                  & Final thesis document & 6 weeks  & M7: Thesis submitted \\
    \bottomrule
  \end{tabular}
\end{table}

\textbf{Key risks and dependencies:}

\begin{itemize}
  \item \textit{Data availability risk:} If real organizational data cannot be obtained,
        synthetic data will be generated from published industry benchmarks.
        (Mitigation: early agreement with partner organization; synthetic fallback planned.)
  \item \textit{Ontology coverage risk:} ESCO may not cover highly specialized CAPEX roles.
        (Mitigation: manual extension nodes linked to ESCO via \texttt{skos:broader}.)
  \item \textit{Integration complexity:} Combining KG embeddings with ML pipelines requires
        tight interface design. (Mitigation: modular architecture with defined APIs between
        components.)
  \item \textit{Expert availability:} Interview scheduling with senior portfolio managers
        may be delayed. (Mitigation: recruit multiple candidates early; parallelize with
        data preparation.)
\end{itemize}

% ─────────────────────────────────────────────────────────────────────────────
\section{Ethical Considerations}
\label{sec:ethics}

If real organizational data is used, appropriate data anonymization and confidentiality agreements
will be established in advance. The recommendation engine is designed to support — not replace —
human decision-making, with explainability mechanisms ensuring transparency of recommendations.
Compliance with GDPR and the EU AI Act is built into the Run phase of the AI-4-SME methodology.
